\chapter{Corusco Tables}
\label{ch:corusco-tables}
\chapterquote{A table column deserves an identity more stable than the integer it happens to occupy today.}{The descriptor rule}

\begin{chaptermap}
Core question & How does Corusco make serious \api{JTable} work less fragile without hiding Swing's real table mechanics?\\
Primary APIs & \api{TableDescriptor}, \api{TableKey}, \api{ColumnKey}, \api{Column}, \api{ColumnDescriptor}, \api{ObservableTableModel}, \api{TableStateController}, \api{TableSelectionBinding}, \api{TableCellValidationBinding}, \api{TableCellTooltipBinding}, \api{TableHeaderTooltipBinding}, and \api{TableHeaderColumnVisibilityMenu}.\\
Swing topics & Model and view coordinates, selection, sorting, filtering, validation, tooltips, persisted layout, hidden columns, editable cells, renderers, and large data.\\
Companion examples & \largeDataExample, \api{ObservableTableModelExample}, \api{TableSelectionBindingExample}, \api{TableStateExample}, \api{GeneratedTableColumnsExample}, \api{TableCellValidationExample}, and \api{TableHeaderTooltipExample}.\\
\end{chaptermap}

\section{A JTable is not a table contract}

\api{JTable} is a powerful Swing component, but it is not the durable contract
of a business table. It knows how to display a \api{TableModel}, manage a
column model, track selection, ask renderers and editors for components,
cooperate with a row sorter, and paint a rectangular view. It does not know
which column identity should survive localization, which row type is expected,
which resource key names a header, which cell value type a column returns, or
how a user's column layout should migrate after a release. Those facts belong
above the component.

\termdef{Table descriptor}{in Corusco} is the Swing-neutral table contract. It
combines a typed \api{TableKey<R>} with ordered executable
\api{Column<R,V>} definitions. Each column has a typed \api{ColumnKey}, a
\api{ColumnDescriptor}, a value extractor, and optionally an updater for
editable cells. A \api{JTable} becomes one presentation of that contract.

\begin{lstlisting}[caption={Descriptor first, Swing table second.}]
TableDescriptor<CustomerRow> descriptor =
        GeneratedCustomerRowTableDescriptor.DESCRIPTOR;

ObservableTableModel<CustomerRow> model =
        ObservableTableModel.of(rows, descriptor);

JTable table = new JTable(model);
\end{lstlisting}

The descriptor owns no rows and no Swing state. It describes meaning. The
observable list owns rows. The Swing table model adapts rows and descriptor to
Swing. The table component owns view state: selection, column model, row
sorter, renderers, editors, focus, keyboard behavior, and painting. This
division is the table version of the core/Swing boundary used throughout the
book.

\begin{figure}[htbp]
\centering
\begin{minipage}[t]{0.48\linewidth}
\includegraphics[width=\linewidth]{\coruscoTableDescriptorFigure}
\end{minipage}\hfill
\begin{minipage}[t]{0.48\linewidth}
\includegraphics[width=\linewidth]{\coruscoTableStateFigure}
\end{minipage}
\caption{Corusco table screenshots show the two halves of the contract: descriptor-backed column meaning and persisted Swing table state by stable ids.}
\label{fig:corusco-table-screens}
\end{figure}

Figure~\ref{fig:corusco-table-screens} gives the chapter's main picture. The
descriptor side names row type, column identity, resources, and value access.
The state side records what users do to the Swing view: widths, order,
visibility, and sorting. A serious table needs both halves. If the descriptor
is missing, user state is stored against weak things such as headers or
indexes. If the view state is missing, a table may be semantically clean but
hostile to users who must resize the same columns every morning.

The distinction is not academic. A table model that hard-codes column indexes
may be shorter on the first day, but it has no single place where column
meaning can be inspected. The header setup knows one set of indexes. The
renderer setup knows another. Validation code knows a third. Preference
persistence may know a localized header. Export may scrape renderer text.
Descriptor-backed tables force these meanings into one typed contract before
Swing is asked to display anything.

\begin{pitfallbox}{the integer that lied}
An integer table coordinate is a position, not an identity. It can answer
"where is this column or row now?" It cannot safely answer "which business fact
is this?"
\end{pitfallbox}

Plain Swing is not the enemy here. A disciplined \api{AbstractTableModel}, a
clear row type, and careful coordinate conversion can produce a good table.
Corusco becomes useful when the same table accumulates several durable
concerns: localization, generated headers, validation targets, row selection,
state persistence, hidden columns, editable records, help topics, export,
testing, or large-data update precision. These concerns exist in plain Swing
too. Corusco names them and gives them typed boundaries.

\section{Rows, columns, and stable identity}

Start with the row. A row type should represent the values the table displays
and the identity the screen needs for selection, refresh, persistence, export,
and master-detail workflow. A row may be an immutable record, a projection
from a domain object, a lightweight DTO, or a mutable model object. What it
should not be is "whatever the current table model happens to return at
column three." The row type is the first step away from index-based thinking.

\begin{lstlisting}[caption={A row type before there is a JTable.}]
public record CustomerRow(
        CustomerId id,
        String displayName,
        CustomerStatus status,
        LocalDate lastOrderDate) {
}
\end{lstlisting}

\api{TableKey<R>} names the table contract for a row type. A screen may have
several tables over the same row type, and each can have a different table
key. A search-results table, an export-preview table, and an audit-history
table may all show customer-like rows but need different columns, persisted
state, and help. The table key is therefore not merely the Java row class. It
is the stable identity of this table contract.

\api{ColumnKey<R,V>} names one column. It carries the row type and the value
type. That type information is not decoration. It prevents a status column key
from being used where a date column is expected, and it lets descriptor,
validation, tooltip, export, and test code talk about the same column without
using a view index. The id string inside the key remains useful for resources,
diagnostics, persistence, and generated source, but application code should
keep the typed key where possible.

\begin{lstlisting}[caption={Typed table and column keys.}]
static final TableKey<CustomerRow> CUSTOMERS =
        TableKey.of("customer/table", CustomerRow.class);

static final ColumnKey<CustomerRow, String> DISPLAY_NAME =
        ColumnKey.of("customer/table/display-name",
                CustomerRow.class, String.class);
\end{lstlisting}

Do not make column ids follow visible headers. A header changes for
localization, copyediting, or density. A column id changes only when the
durable identity changes. \api{"customer/table/status"} can keep its identity
when the English header changes from \api{"State"} to \api{"Status"}. If a
column is split into two different business facts, then a new id and a table
state migration may be justified. The id policy should make that distinction
visible.

The same row value may have several identities in a screen. The row object may
be equal by value. The domain object may have a \api{CustomerId}. The table
model row index changes when rows are inserted or removed. The view row index
changes when sorting or filtering changes. The selected row object may become
stale after refresh. A robust table design decides which identity is used for
which job. Persisted selection across refresh normally wants a domain id.
Immediate selection binding may publish a row value. Swing repaint wants a
model row. The visible table selection wants a view row.

\begin{figure}[htbp]
\centering
\includegraphics[width=0.88\linewidth]{\tableCoordinatesFigure}
\caption{A table has several coordinate systems. Corusco descriptors and bindings keep durable row and column identity separate from current Swing view positions.}
\label{fig:corusco-table-coordinates}
\end{figure}

Figure~\ref{fig:corusco-table-coordinates} belongs near the beginning because
many table bugs are not renderer bugs or model bugs. They are coordinate
bugs. A user selects the third visible row, but the presenter saves model row
three after sorting. A tooltip asks for view column two, but validation stored
model column two. A table state store writes the English header instead of the
column persistence id. Once the coordinate systems are named, the fixes become
straightforward.

Stable identity also improves documentation. A table descriptor can be cited in
the book, a test, a migration note, or a resource audit. A phrase such as
"the customer status column" can correspond to a \api{ColumnKey}. Without that
key, the phrase is implemented as a header string in one place, an index in
another, and a renderer branch in a third. The human concept remains singular;
the code has split it into accidents.

\section{Descriptor anatomy}

\api{ColumnDescriptor<R,V>} contains the metadata side of a column. It has a
column key, header resource key, optional tooltip resource key, optional help
topic, persistence metadata, default presentation state, and declared
capabilities. It does not read row values. It does not update rows. It does
not know Swing components. It says what the column is and what default
presentation policy belongs to it.

\api{Column<R,V>} is the executable side of a column. It wraps the descriptor
with a getter and, for editable columns, an updater. The getter reads a value
from a row. The updater returns a replacement row after an edit. This design
fits immutable records naturally, but it also works with mutable row objects
when the updater returns the mutated object. The important rule is that the
adapter receives a coherent row value after editing.

\begin{lstlisting}[caption={A handwritten descriptor-backed column.}]
ColumnDescriptor<CustomerRow, String> nameDescriptor =
        new ColumnDescriptor<>(
                DISPLAY_NAME,
                CustomerResources.DISPLAY_NAME_HEADER,
                CustomerResources.DISPLAY_NAME_TOOLTIP,
                CustomerHelp.DISPLAY_NAME,
                new ColumnDefaults(220, true, 0),
                ColumnCapabilities.readOnly());

Column<CustomerRow, String> name =
        Column.readOnly(nameDescriptor, CustomerRow::displayName);
\end{lstlisting}

Generated \api{@SwingTable} companions write this kind of code mechanically.
For a record component annotated with \api{@Column}, generated code can create
the column key, resource keys, descriptor, getter, and table descriptor. The
generated form is not magic; it is typed Java that a careful developer might
write by hand and would not want to repeat for every table.

\api{ColumnDefaults} and \api{ColumnPersistence} deserve attention. Defaults
describe initial view policy: width, visibility, and default order. Persistence
describes how user state should be stored and bounded: persistence id, default
width, and maximum restored width. Defaults answer "what should the table look
like when no user state exists?" Persistence answers "how do saved values map
back to this column?" These are related but not identical.

Capabilities describe what the screen may offer. A column may be sortable,
filterable, hideable, or editable. This is not only for generated menus. A
filter builder should not show columns marked non-filterable unless
application policy overrides the descriptor. A column visibility menu should
respect hideability. A generated editable column should have an updater. A
renderer may still render any column, but user-facing operations should not
pretend all columns support every capability.

The visible header is not the descriptor. The descriptor points to a resource
key. A Czech translation, an English copy edit, or a shorter header for a
narrow screen should not change validation targets or persisted table state.
Resource lookup belongs near presentation, where missing resources can be
handled by policy and where localized strings can be applied to headers,
menus, tooltips, and accessibility.

\begin{lstlisting}[caption={Using descriptor metadata at the Swing edge.}]
Column<CustomerRow, ?> column = descriptor.column(modelColumn);

String header =
        resources.require(column.descriptor().headerKey());
String tooltip =
        resources.resolve(column.descriptor().tooltipKey(), "");
\end{lstlisting}

Help topics also belong in descriptors. A table cell is often visually small,
but the column may have a real business explanation. F1 on a header or cell can
route to the column's help topic if the descriptor supplies one. This is much
stronger than guessing help from the current header text. It also lets a
command palette, documentation generator, or accessibility adapter discover
help metadata without walking Swing components.

\begin{coruscobox}{what the descriptor buys}
The descriptor gives one durable place for row type, column identity, resource
keys, help topics, value extraction, editing behavior, default widths,
persistence ids, capabilities, generated tests, and documentation references.
Swing remains responsible for rendering and interaction, but it does not have
to invent table meaning from indexes and header text.
\end{coruscobox}

\section{Observable models and row changes}

\api{ObservableTableModel<R>} is a Swing \api{TableModel} backed by an
\api{ObservableList<R>} and a \api{TableDescriptor<R>}. It is EDT-confined.
Construct it on the EDT. Mutate the source list on the EDT while the model is
subscribed, or wrap the source with \api{EdtObservableList} before constructing
the model. Without an explicit dispatcher, off-EDT source changes should fail
fast instead of firing Swing events on the wrong thread.

\begin{lstlisting}[caption={Observable row source to JTable.}]
ObservableList<CustomerRow> rows = presenter.rows();
EdtObservableList<CustomerRow> swingRows =
        EdtObservableList.of(rows);

ObservableTableModel<CustomerRow> model =
        ObservableTableModel.of(swingRows,
                GeneratedCustomerRowTableDescriptor.DESCRIPTOR);

table.setModel(model);
\end{lstlisting}

The table model answers Swing's questions by consulting the descriptor. Row
count comes from the observable list. Column count comes from descriptor
columns. Column class comes from the column value type. Column names default
to header key ids until resources are applied at a higher layer. Cell values
come from the column getter. Editability comes from column update capability.
The model is small because row and column meaning are elsewhere.

The model subscribes to the row source and translates list changes into table
events. An inserted range becomes \api{fireTableRowsInserted}. A removed range
becomes \api{fireTableRowsDeleted}. A replacement becomes
\api{fireTableRowsUpdated}. A clear becomes a delete range. This is why row
source precision matters. The model cannot fire precise Swing events if the
row source reports every small edit as "clear and refill everything."

\begin{coruscobox}{list precision becomes table precision}
Corusco tables are not faster because they ignore Swing. They are faster and
more stable when observable row sources report precise changes that Swing can
translate into precise table events.
\end{coruscobox}

Large data makes this rule visible. A table with 200 rows can survive broad
refreshes. A table with 200,000 rows cannot casually rebuild itself on every
edit. Precise row replacement preserves scroll position, selection, sorter
state, renderer caches, and user confidence. Broad changes are sometimes
correct, especially after a new search query or refresh from a different data
source. They should be broad because the workflow is broad, not because the
model lacks vocabulary for smaller changes.

\begin{figure}[htbp]
\centering
\includegraphics[width=0.88\linewidth]{\largeDataFigure}
\caption{Large table examples make row-event precision visible. Descriptor-backed columns help, but row-source update policy still decides how much of the table Swing must revisit.}
\label{fig:corusco-table-large-data}
\end{figure}

Figure~\ref{fig:corusco-table-large-data} is intentionally not a glamorous
table. It is the kind of dense data screen where table discipline matters. A
cell renderer that allocates too much, a model that fires full refreshes, or a
selection policy based on view rows will show up quickly. Corusco's table
model gives the screen a typed contract, but performance still depends on
using Swing's event model accurately.

Editable cells cross from Swing editor to row source. \api{ObservableTableModel}
checks whether the column is editable, calls the column updater, replaces the
row in the source list, suppresses the duplicate source event for its own
immediate update path, and fires a cell update. The updater receives the
current row and the new value cast to the column's declared value type. For
record rows, the updater normally calls the canonical constructor or a wither.
For mutable rows, it may mutate and return the same object. In both cases, the
row source must represent the semantic edit.

Do not put all parsing and validation inside the table editor. The editor may
collect text. The row model, form model, or validation model should own
semantic validity. A cell validation binding can then display structured
problems. An editor that rejects every intermediate invalid keystroke may feel
hostile; an editor that commits invalid semantic data into the row may corrupt
state. Treat editing as a small workflow, not only a call to \api{setValueAt}.

Closing the model matters. \api{ObservableTableModel} implements the binding
lifecycle because it owns a subscription to the row source. When the view
closes, close the model or the binding scope that owns it. Otherwise a row
source can keep notifying a table model after the table is gone. That leak may
not appear in a short demo, but it appears in long-running desktop sessions.

\section{Selection and coordinate systems}

Swing selection lives in view coordinates. Presenter state should not. A table
may be sorted or filtered, so the selected view row must be converted to a
model row before it becomes application state. If the table is refreshed, the
old model row may not name the same domain object. If a row sorter is active,
the visible row order is not the source order. These are ordinary Swing facts,
and a serious table must name them rather than pretend they are rare.

\api{TableSelectionBinding} binds a \api{JTable} selection to a writable
selected model-row value and optionally a selected row value. User-originated
selection writes \api{ChangeOrigin.USER}. Presenter-originated changes to the
model-row value select the corresponding view row. The binding uses Swing's
view/model conversion APIs and listens to selection, table-model, and row-sorter
changes so selection remains coherent after sorting or structural updates.

\begin{lstlisting}[caption={Selection as model state.}]
WritableValue<Integer> selectedModelRow = SimpleValue.empty();
WritableValue<CustomerRow> selectedRow = SimpleValue.empty();

Binding selection = TableSelectionBinding.bind(
        table, model, selectedModelRow, selectedRow);
\end{lstlisting}

This binding is intentionally single-row oriented. Swing can support multiple
selected rows, intervals, lead selection, and anchor selection. If a screen
needs multi-selection, model that explicitly with a selected-id set or selected
row collection. Do not pretend that a single lead row represents a batch
operation. Delete selected rows, export selected rows, and compare selected
rows deserve a multi-selection model.

Configure any row sorter before creating the binding. If the table receives a
replacement sorter, close and recreate the binding. This rule follows from
ownership. The binding attaches to the current sorter to refresh selection
after sort changes. If another part of the view silently swaps sorters, it has
changed a relationship the binding owns. Recreate the relationship rather than
hoping listener state still matches the table.

Filtering raises a design question: does filtering belong in Swing's
\api{RowFilter} or in a presenter-owned filtered list? If filtering is purely a
view choice, such as "hide closed rows in this table only", a Swing sorter may
be enough. If filtering affects command state, status text, export, summary
counts, master-detail selection, or other screen areas, put the filter in the
presenter or row source and let the table display the result. A filter that
changes application meaning should not live only inside a \api{JTable}.

Selection preservation after refresh is policy, not an accident. If rows are
reloaded from a service, should the previously selected customer remain
selected if it still exists? Should selection clear if the row disappears?
Should the nearest visible row become selected? The answer depends on the
workflow. The implementation should use durable row identity where continuity
matters. A model row index from before refresh is rarely durable enough.

The testing consequence is direct. A weak table test says "select row two and
assert row two is selected." A stronger test says "select the row whose id is
42, sort by name, refresh with the same id present, and assert the selected row
identity is still 42." The second test follows the user's mental model. The
user selected Grace Hopper, not whatever happened to occupy view row seven.

\section{State, visibility, and migration}

\termdef{Table state}{in Corusco} includes column order, widths, visibility,
and sort state. Users expect serious business tables to remember these
choices. Persisting them by localized header text is a mistake. Persisting
them by stable table ids and column persistence ids is the right shape. A
header may change for translation; a column persistence id should remain
stable unless the column's durable meaning changes.

\api{TableStateController} is the bridge between toolkit-neutral table
metadata in \api{TableDescriptor}/\api{TableState} and Swing's mutable
\api{TableColumnModel}. It reads stable column persistence ids from the
\api{ObservableTableModel} descriptor, merges saved state with current
descriptor defaults, applies column order, visibility, width bounds, and sort
keys to the table, then listens for Swing column and sorter changes. It owns
the listeners and save scheduler. It does not own the table, model,
descriptor, or store.

\begin{lstlisting}[caption={Installing persisted table state.}]
TableStateStore store =
        new PreferencesTableStateStore(preferences);

TableStateController<CustomerRow> controller =
        TableStateController.install(table, model, store);
\end{lstlisting}

State writes are debounced on the EDT so resize and drag gestures coalesce
into one store write. Programmatic changes such as column visibility changes
also schedule persistence. Closing the controller removes Swing listeners,
writes pending state, and flushes the store. After close, mutation and capture
methods should reject use. That strict lifecycle is useful because table state
controllers otherwise become quiet leak sources.

Hiding a column is not the same as removing it from the model. A hidden column
may still exist in the descriptor and row model. It may still participate in
export, validation, help, or future visibility. The table state controller can
hide and restore columns by persistence id while keeping the descriptor
stable. A visibility menu should list columns by descriptor metadata and
localized resources, not by current header text scraped from the table.

\api{TableHeaderColumnVisibilityMenu} is useful for this reason. It can present
a popup trigger for column visibility while using descriptor identities under
the surface. The user sees localized names. The table state receives stable
ids. The code avoids a common bug where hiding a column removes it from the
model and shifts every remaining index.

Persisted table state outlives one application version. Columns are renamed,
split, removed, or added. \api{TableStateMigration} lets loaded state be
adjusted before descriptor defaults are merged. Use migration sparingly and
document it. Table ids and column ids should be stable by default. Migration
is for real schema changes, not for casual id renaming. A migration that
exists only because a header was copyedited is a sign that identity and
presentation were confused.

Corusco separates table state from the place where state is stored.
\api{InMemoryTableStateStore} is useful for tests and examples.
\api{PreferencesTableStateStore} is useful for desktop applications that want
per-user settings through Java preferences. Larger applications may provide a
store backed by workspace files, application configuration, or a database. The
store should not need Swing. It stores \api{TableState} by stable ids. The
controller knows Swing. The descriptor knows table meaning.

\begin{tabularx}{0.96\linewidth}{@{}p{0.31\linewidth}X@{}}
\toprule
Observation & Review question\\
\midrule
Column id changed with label text & Was a resource edited when an id should have stayed stable?\\
Selection breaks after sorting & Is the code mixing view row and model row?\\
Persisted widths reset after release & Did table or column persistence ids change?\\
Renderer runs validation & Should validation problems be computed before rendering?\\
Hidden column removed from model & Is visibility being confused with descriptor membership?\\
Refresh clears selection & Are row identity and precise list events preserved?\\
\bottomrule
\end{tabularx}

\section{Validation, tooltips, help, and accessibility}

Table validation problems need stable targets. A problem at view row 4,
column 2 is not stable under sorting. A problem for a row identity and
\api{ColumnKey} is. \api{TableCellValidationBinding} decorates cells by
finding problems for the row and column being rendered. The problem model
remains independent of current view coordinates. Rendering displays an already
computed fact; it should not run the validation rule.

\begin{lstlisting}[caption={Cell validation observes structured problems.}]
ReadableValue<ProblemSet> problems =
        presenter.tableProblems();

Binding validation =
        TableCellValidationBinding.bind(table, model, problems);
\end{lstlisting}

The row side of a problem target requires a screen decision. Some screens can
target by row object because rows are stable values for the lifetime of the
table. Some should target by domain id because rows are replaced after
refresh. Some validation is transient and can target the current model row
only. The column side should still use \api{ColumnKey}. The important point is
that validation identity is explicit rather than inferred from current view
coordinates.

Tooltips are another shared surface. Header tooltips may come from column
descriptors and resources. Cell tooltips may come from validation problems,
static column help, formatted values, disabled reasons, or help availability.
Corusco provides table tooltip bindings so tooltip ownership is not scattered
through renderers and header mouse listeners. The same composition principle
from commands and forms applies: several sources contribute to one tooltip
slot; a policy composes them.

\begin{lstlisting}[caption={Header tooltip binding uses descriptor resources.}]
Binding headerTooltips =
        TableHeaderTooltipBinding.bind(table, model, resources);
\end{lstlisting}

Cell tooltips should remain cheap. A tooltip query may happen frequently as
the mouse moves. It should not call services, run full validation, or allocate
large formatting structures. It should look up descriptor metadata, current
row value, and already computed problems. If a cell requires expensive
explanation, precompute the explanation or route to help instead of doing the
work on hover.

Help topics make table documentation precise. F1 on a table header or focused
cell can route to a column help topic. A help request may include the column
key, row identity, current value, or validation problem context. It should not
derive the topic from the visible header. Header text is presentation. The
help topic is durable routing metadata.

Accessibility follows the same metadata. A header resource can contribute to
the accessible column name. A tooltip resource or help topic can contribute to
an accessible description. A validation problem can be exposed through a
problem summary, tooltip, or cell renderer state. A color-only cell marker is
not sufficient for all users. The descriptor gives the screen a typed place to
attach accessible meaning; the renderer and table bindings decide how to
present it.

Validation and tooltip bindings should be closed with the view. They install
renderers, mouse listeners, or table integrations that belong to a particular
\api{JTable}. If a table model outlives the view, the old view should not
continue to receive problem updates. The same lifecycle rule appears
throughout Corusco Swing: install relationships through scopes and close them
when the owner closes.

\section{Editing cells without losing the row}

Editable cells look simple until the first invalid intermediate value appears.
Swing editors are components. They can hold text or selection state that is
not yet committed. The table model row is application data. The two should not
be confused. Commit should happen through Swing's editor lifecycle, and the
column updater should produce a coherent row update only when the edit is
ready to enter the row source.

For immutable rows, an editable column normally returns a replacement row:
\api{new CustomerRow(row.id(), value, row.status())} or a generated wither-like
constructor call. The observable list then replaces the row and fires a precise
row update. Selection by row object may need to be refreshed; selection by
domain id is usually more stable. For mutable rows, the updater may set a
property and return the same row, but the observable list still needs to fire a
change event so Swing repaints and dependent state updates.

The value type in \api{ColumnKey<R,V>} is useful during editing because
\api{Column.update} casts the supplied value to the declared type before
calling the updater. This catches accidental editor/model mismatches close to
the table boundary. It does not replace validation. A \api{BigDecimal} value
of \api{-1} may be type-correct and semantically invalid. Type, parse, and
validation are separate concepts.

Some tables should not commit edits directly to the domain row. A price table
may need a transactional edit model. An invoice-line table may need row-level
validation before commit. A remote grid may need optimistic updates and
rollback. Corusco's descriptor model still helps by naming columns and row
identity, but the editable workflow may need a form model per row, an edit
buffer, or a command that applies staged changes. Do not force every table
edit through \api{setValueAt} merely because Swing provides that method.

Active-editor commit matters for commands. If the user is editing a cell and
presses Save, should the active editor commit first? Usually yes. If the
editor cannot commit because the value is invalid, Save should not proceed as
though the old row value were current. Dialog chapters discuss active-editor
commit for forms; tables need the same discipline. The command handler should
know whether pending cell edits have been committed, cancelled, or rejected.

Cell editors also affect focus and keyboard behavior. Enter may commit the
cell, move to the next row, or invoke the default button depending on table
policy. Escape may cancel the editor or close a dialog. These policies should
be reviewed explicitly in dense business screens. Descriptor-backed columns
do not decide key behavior, but they let editor installation and command
policy know which column is being edited.

\section{Rendering and scale}

Swing renderers are already optimized in one important way: one renderer
instance stamps many cells. Most performance problems come from violating that
model by allocating expensive objects, doing business work, or retaining row
state inside renderer components. A renderer cost that seems trivial for one
cell is paid for many visible cells and repeated repaints. Keep validation,
service calls, filtering, and remote lookups out of renderers.

\begin{pitfallbox}{renderer work is multiplied}
A renderer cost that seems trivial for one cell is paid for many visible cells
and repeated repaints. Keep validation, service calls, and row filtering out
of renderers.
\end{pitfallbox}

The descriptor does not make rendering disappear. It makes rendering less
lonely. A renderer can ask what column it is painting through a stable
\api{ColumnKey}; it can read display policy from descriptor metadata; it can
cooperate with table validation and tooltip bindings that use the same column
identity. This is different from a renderer that receives an integer column,
switches on that integer, and silently assumes that no one will move, hide, or
insert a column later.

Start with ordinary Swing renderers. A descriptor-backed table model and
selection binding are enough for many screens. Introduce optimized renderers
only when there is a measured reason or a known dense workload. A timestamp
column that repaints thousands of cells per second may deserve a cached
formatter and a renderer installed by column key. A boolean status column may
deserve a finite-state renderer. A risk column may deserve a custom badge.
Each of these is still Swing presentation. It should not compute validation,
call services, or rewrite row identity.

Look-and-Feel discipline matters. A descriptor can say that a column
represents severity. It should not force one hard-coded red into every theme.
Renderers should prefer table foreground, selection foreground, disabled
foreground, UI defaults, or application theme keys. A renderer that is fast
but unreadable in dark mode is not finished. A renderer that hides state
behind color alone is not finished either; the same descriptor identity should
lead to tooltip, accessibility, export, and testable problem state.

Large data changes the renderer review. A table with hundreds of visible rows
and frequent updates needs renderers that reuse formatters, avoid per-cell
allocation, and do not trigger layout changes. Icons should be cached. Date
and number formatters should be prepared outside the hot path when possible.
Tooltips should avoid recomputing expensive strings on every mouse move.
Selection and row updates should preserve enough state that the table does
not repaint everything unnecessarily.

Sorting and filtering also affect perceived performance. Swing's
\api{TableRowSorter} can be useful for moderate row counts and local data. For
large remote data, sorting and filtering may belong in the service or
presenter-owned row source. The descriptor still names columns and
capabilities. The application decides where the operation runs. Do not let a
small demo sorter become production policy for a million-row query.

Printing, export, and copy are rendering-adjacent but not renderer work. They
should use descriptor metadata and row values, not scrape painted components.
The display renderer may show an icon for status; export may need the status
code or localized text. A descriptor gives the export path the same column
identity without depending on how the current Look-and-Feel paints the cell.

\section{Generated and handwritten tables}

\api{@SwingTable} generation is the shortest path for ordinary record-backed
tables. The annotated record declares stable table and column facts. The
processor emits column keys, resource keys, descriptors, a table descriptor,
and binding helpers. Generated columns call record accessors. Editable
generated columns can produce replacement records through constructor calls
when supported. The generated code is ordinary Java source and should be read
when behavior is surprising.

\begin{lstlisting}[caption={Annotated table row source.}]
@SwingTable(id = "customer/search")
public record CustomerRow(
        @Column(width = 180)
        String name,

        @Column(width = 80, sortable = false)
        int orders) {
}
\end{lstlisting}

Generated tables are not a layout system. They do not decide where the
\api{JTable} belongs in a MigLayout panel, how tall the table should be, which
commands surround it, or whether it appears in a split pane. They generate the
metadata and adapters that make the table meaning stable. The view still owns
Swing composition.

Handwritten descriptors remain useful. A GIS layer table, paint-editor object
table, log-event stream, or diagnostic table may have dynamic columns,
specialized renderers, or column definitions that do not fit record
annotations. Handwriting a \api{TableDescriptor} is not a failure. It is the
same contract expressed directly. The rest of Corusco's table stack can still
use it: observable model, selection binding, state controller, tooltips,
validation, and tests.

The migration path from plain Swing is incremental. First define row type and
column keys. Then replace index switches with columns and descriptors. Then
adapt the existing row list to an observable list. Then install
\api{ObservableTableModel}. Then move selection state out of table listeners.
Then add table state by stable ids. Then add tooltip, validation, export, and
generated resources as needed. Each step should reduce one hidden convention.

\begin{migrationbox}{from JTable conventions to descriptors}
Do not migrate by rewriting every table at once. Migrate one table contract:
row type, column keys, descriptor, observable model, selection binding, and
state persistence. Keep the visible table recognizable while the internal
identity becomes stable.
\end{migrationbox}

For a small diagnostic table, plain Swing may still be the right tool. The
threshold is not row count alone. A table becomes a Corusco table when column
identity matters outside immediate painting: persisted widths, generated
headers, localization, validation targets, row selection, column visibility,
tooltips, tests, editable record rows, or generated documentation. If none of
those apply, the descriptor may be ceremony. If several apply, the descriptor
is usually cheaper than a network of conventions.

The companion examples show several slices. \api{ObservableTableModelExample}
focuses on rows and descriptor-backed values. \api{TableSelectionBindingExample}
focuses on sorting and selection conversion. \api{TableStateExample} focuses
on persisted state and migration. \api{GeneratedTableColumnsExample} shows the
generated path. \api{TableCellValidationExample} and
\api{TableHeaderTooltipExample} show feedback surfaces. Read them as small
chapters in table ownership rather than as unrelated demos.

\section{A table lifecycle in one screen}

It is helpful to trace one realistic table from construction to teardown. The
screen begins before a \api{JTable} exists. The source declaration or
handwritten descriptor defines the row type and column identities. At this
stage there are no view rows, no table columns, no selected index, and no
renderer calls. There is only a contract: this screen has a customer table;
it has columns such as name, state, and last order; each column has stable
identity and resource keys; some columns may be editable; table layout should
be persisted by ids that survive localization.

The presenter then creates or receives an observable row source. This source
is not Swing state. It is the current set of rows the screen is responsible
for. It may be loaded from a repository, filtered by presenter logic, sorted
by a business rule, mapped from a larger model, or adapted from another
event-list library. The important point is that changes are stated as changes.
One row replacement should not become a declaration that the whole table
universe collapsed. A batch load may still be broad, but it should be broad
because the workflow is broad.

Only after descriptor and rows exist does the Swing table model appear. The
\api{ObservableTableModel} adapts row source and descriptor to Swing's
\api{TableModel} protocol. Swing asks row count, column count, column class,
column name, cell value, editability, and value update. The model answers
through descriptor metadata and row values. When the source list changes, the
table model translates list change sets into table events. The \api{JTable}
does not need to know where rows came from. It only sees a model that obeys
Swing's contract.

The table component is then configured for presentation. The view may install
a row sorter, custom renderers, column width defaults, selection mode, grid
policy, row height, popup menu, and keyboard actions. These are real Swing
decisions and should remain visible in view code. Corusco should not hide them
behind a generated GUI builder. What Corusco supplies is stable metadata that
makes those decisions less brittle. A renderer can be installed for the state
column by key. A popup menu can operate on selected row identity. A header
tooltip can resolve a descriptor resource key.

Selection is the next boundary. A user selects a view row. If sorting or
filtering is active, that row is not the model row. If the table has been
refreshed, the selected index may not name the same business object it named a
moment ago. The selection binding performs mechanical Swing conversion and
publishes model-row or row-value state according to the screen's needs. The
presenter reacts to this state, not to the table's transient selected view
row. This keeps master-detail loading, command enablement, and status text on
the same side of the boundary.

Validation follows the same identity rule. A problem belongs to a row and a
column. A renderer or tooltip binding can ask whether the cell currently being
painted has a problem, but it should not run the rule that discovers the
problem. The rule belongs to form, row, presenter, or validation model code.
The table presentation displays an already computed structured fact.

Persistence enters after the table's columns exist. The state controller
restores width, order, visibility, and sort state from stable ids. It then
listens to column and sorter changes and writes snapshots through a store.
This is view state, not table meaning. A user may decide that Status should be
the first visible column. The descriptor still knows that Status is the status
column. If the descriptor id changes casually, old user preferences, tests,
validation, export, and help all become suspect.

Refresh tests the whole design. A background task loads new rows. The result
returns on the EDT or through an EDT-adapted list. The row source applies a
precise update where possible. Selection is preserved by identity when the
user expects continuity. Table state remains attached to column ids.
Validation and tooltip bindings observe the new problem state. Renderers
repaint visible cells. If any part of this chain was based on localized
headers, stale view indexes, or hidden listeners, refresh is where the bug
usually appears.

Finally, teardown must be ordinary. The view closes the table model
subscription, selection binding, table state controller, tooltip bindings,
validation bindings, and any renderer installer handles that promised
restoration. If the row source changes after the view closes, the old table
should not receive events. If a table state save was scheduled, it should
either finish under its owner or be cancelled by policy. A table whose
lifecycle is unclear will eventually show duplicate listeners, stale selection
writes, or preference updates from a screen the user has already closed.

\begin{lstlisting}[caption={One reviewable table assembly outline.}]
ObservableList<CustomerRow> rows = presenter.rows();
TableDescriptor<CustomerRow> descriptor =
        GeneratedCustomerRowTableDescriptor.DESCRIPTOR;

EdtObservableList<CustomerRow> swingRows =
        EdtObservableList.of(rows);
ObservableTableModel<CustomerRow> model =
        ObservableTableModel.of(swingRows, descriptor);

JTable table = new JTable(model);
table.setRowSorter(new TableRowSorter<>(model));

scope.add(TableSelectionBinding.bind(
        table, model,
        presenter.selectedModelRow(),
        presenter.selectedRow()));

scope.add(TableCellValidationBinding.bind(
        table, model,
        presenter.tableProblems()));

scope.add(TableHeaderTooltipBinding.bind(
        table, model,
        resources));

scope.add(TableStateController.install(
        table, model,
        tableStateStore));
\end{lstlisting}

The outline is not a recipe for every table. Some tables do not need
validation. Some do not need persisted state. Some use a presenter-owned
sorted list rather than a Swing row sorter. Some have multi-selection. The
value of the outline is that every installed relationship has a name and a
scope. Remove one concern only when the screen truly does not need it, not
because the concern was hidden inside a listener.

\section{Large data and refresh workflows}

Large data does not begin at a magic row count. It begins when table updates
become visible to users: a refresh moves the scroll position, a sorter stalls,
a renderer allocates too much, selection disappears, or a background result
arrives after the user has moved on. The large-data chapter studies lists in
general. In a descriptor-backed table the same concerns appear with Swing
specific consequences. The row source must report meaningful changes. The
table model must translate them on the EDT. Selection must be preserved by
identity when policy says so. Renderers must stay cheap. Table state must not
be reset merely because rows changed.

There are three common row-source shapes. The first is a small local list
loaded all at once. It can use \api{ObservableArrayList} directly on the EDT.
The second is a medium local list whose changes may be computed away from the
EDT and delivered through \api{EdtObservableList}. The third is a large or
remote source where the visible list is only a window, filter result, or query
result. The descriptor works in all three shapes, but the update policy
differs. A small list can be simple. A large list needs refresh generations,
precise replacement, and sometimes service-side sorting or filtering.

A refresh workflow should declare what happens to old rows while new rows are
loading. Keeping old rows visible is often best because the user can still
read the table, but it requires a visible stale or loading state so the user
knows a refresh is in progress. Clearing rows immediately may be appropriate
when old rows would be misleading, but it destroys selection and scroll
context. Replacing rows only after a successful load preserves context on
failure. None of these choices belongs in \api{JTable}; they belong to the
presenter and row source.

When rows arrive, do not assume that \api{clear} plus \api{addAll} is the only
honest update. If the service returns stable row ids, the presenter can compute
a diff: removed ids, inserted ids, moved rows, and replaced row values. For
small tables the diff may not be worth the code. For large tables or
frequently refreshed tables it can preserve selection, sorter state, and
rendering work. The descriptor does not compute the diff; it makes the table
model able to benefit from a precise diff when the row source supplies one.

Selection preservation is usually a row-identity problem. Suppose the user
selects customer 42, then refresh runs. If customer 42 is still present, the
screen may preserve the selection even if the row moved. If customer 42 is
gone, the screen may clear detail or select the nearest visible row. If the
table is sorted by name, preserving "model row five" is not the same as
preserving customer 42. The selected-row value can help immediate workflow;
a selected-id value often helps refresh continuity.

\begin{lstlisting}[caption={Refresh by row identity, not by old view row.}]
CustomerId previouslySelected = selectedCustomerId.value();
rows.replaceAll(diffedRowsFrom(serviceRows));

if (previouslySelected != null) {
    rows.indexOfFirst(row -> row.id().equals(previouslySelected))
            .ifPresentOrElse(
                    selectedModelRow::setValue,
                    () -> selectedModelRow.setValue(null));
}
\end{lstlisting}

The exact list API may differ, but the policy is the point. The selected fact
is a customer id. The table row after refresh is found by that id. Swing view
coordinates are recalculated after the model changes. This is slower to write
than \api{table.getSelectedRow()} and much cheaper to maintain after sorting,
filtering, and refresh appear.

Sorting location matters. Swing's \api{TableRowSorter} is convenient and works
well for many local tables. It sorts the view, leaving the model order
unchanged. A presenter-owned sorted list changes the model order before Swing
sees it. A service-side sort returns data already sorted. Each option has a
different owner. Use Swing sorting for local view preference. Use presenter or
service sorting when sorted order affects other screen state, paging, remote
query cost, or exported order. The descriptor can mark which columns are
sortable; it does not decide where sorting runs.

Filtering has the same three possible owners. A Swing row filter is local and
view-oriented. A Corusco filtered list can feed several views or command
states. A service filter can reduce remote data volume. Do not choose the
owner only because one API is shorter. Ask who needs to know the filtered set.
If a status label says "34 matching customers", an export command exports the
filtered set, and a detail panel reacts to filtered selection, the filter is
probably presenter state rather than only table view state.

Virtualized data requires an even stronger boundary. A Swing \api{JTable}
still asks for row count and cell values, but the row source may represent a
window into a remote result. Descriptors remain useful because columns have
stable identities and value extractors for loaded rows. Selection should be by
domain id. Table state should be independent from current page. Renderers must
handle loading placeholders or absent values deliberately. A virtual table is
not just a bigger local list; it is a table whose row source has loading
semantics.

Glazed Lists interop fits naturally when an existing application already uses
\api{EventList}. The adapter boundary should still be explicit. Let Glazed
Lists own its event list and advanced filtering if that is the established
tool. Adapt it to Corusco's observable-list contract at the Swing boundary
where needed. Do not mix two event systems in the same method. The table model
should see one observable row source with clear EDT delivery.

Large table diagnostics should print stable ids. A log entry saying
"full refresh of customer/table, 80,000 rows, previous selection customer 42
not found" is useful. A log entry saying "selected row -1" is not. Include
table key, row counts, changed ranges, sort keys, filter descriptions, and
column ids when diagnosing performance or selection problems. Typed table
identity makes such diagnostics natural.

\section{Commands, export, and surrounding screens}

A table rarely stands alone. It sits in a screen with commands: Refresh,
Delete, Edit, Export, Copy, Toggle Column, Show Details, Open Help, and maybe
Select All. These commands should not scrape the \api{JTable} for meaning.
They should use descriptor identity, selected-row state, row ids, and command
enablement derived from presenter state. The table component is the user's
surface; it should not be the command model.

Delete is the simplest example. A delete command should be enabled when a
deletable selection exists and the user has permission. The command should
receive selected row identity from presenter state. It should not call
\api{getSelectedRow} and convert coordinates in the handler. If the table is
sorted and selection binding already publishes selected id or row, the command
can ignore the current view coordinate. This makes toolbar, menu, popup, and
keyboard paths consistent.

Popup menus are a table-specific trap. A right-click may happen on a row that
is not currently selected. Some applications select the row under the pointer
before showing the popup. Others keep selection unchanged and offer commands
for the clicked row. Either policy can be right, but it should be explicit.
Use row identity and selection binding to implement the policy. Do not let a
popup action silently operate on whatever row happened to be selected before
the mouse press.

Export should be descriptor-oriented. A visible table may hide columns,
reorder columns, show icons, use localized headers, and format values for
density. An export may need all columns, visible columns only, a product-owned
export order, raw values, localized display values, or separate export labels.
The \api{JTable} cannot infer this product decision. The descriptor can supply
stable column identity and value extraction; an export descriptor or export
policy can decide which columns and value representations are exported.

\begin{lstlisting}[caption={Export starts from descriptors and rows, not renderers.}]
List<Column<CustomerRow, ?>> columns =
        exportPolicy.columnsFor(descriptor, tableState);

for (CustomerRow row : selectedRows.value()) {
    for (Column<CustomerRow, ?> column : columns) {
        export.writeCell(column.key(), column.value(row));
    }
}
\end{lstlisting}

Copy to clipboard is similar but more presentation-oriented. Copying visible
cells may reasonably follow current view order and visible columns. It still
should not scrape renderer components. Convert selected view rows to model
rows, read row values through descriptor columns, resolve headers through
resources, and format values through a copy policy. The visible table state
contributes ordering and visibility; the descriptor contributes meaning.

Column visibility commands are table commands too. A column chooser or header
popup toggles visibility by column persistence id. It should honor
\api{hideable} capability and update the \api{TableStateController}. The
command text comes from resources. The selected state comes from current table
state. Tests can assert that toggling \api{customer/table/status} hides the
status column, not that "the third visible column" disappeared.

Master-detail screens should treat the table as a selection source, not as
the detail owner. Selecting a row publishes selected identity. The presenter
loads detail or creates a form model. Commands observe detail state. The table
may show validation markers or stale-row markers, but it should not own the
detail form. This separation keeps the table usable even when detail loading
is asynchronous, fails, or is cancelled by a new selection.

Navigation and help also use descriptor identity. Double-click on a row may
open detail. Enter may activate the primary row command. F1 may open help for
the focused column. A context menu may show column help. These gestures are
Swing input events, but their meaning comes from row and column identity.
Generated or handwritten descriptors provide the stable terms that command
and help systems can use.

Accessibility review should include surrounding commands. A table with icon
status cells, hidden columns, and popup-only operations can be difficult for
keyboard and assistive-technology users. Descriptor metadata can feed header
names, tooltip text, help topics, and command labels. The view still needs
keyboard access, focus traversal, and accessible descriptions. Do not assume a
well-typed descriptor automatically creates an accessible table.

The table's status text should be model-driven. "125 customers", "34 matching",
"Loading", "Stale results", "3 selected", and "Selection hidden by filter" are
not table paint details. They are facts derived from row source, filter state,
selection state, task state, and table state. A status label bound to those
facts is easier to test than a series of \api{setText} calls inside listeners.

\section{Migration from legacy JTable code}

Legacy Swing tables usually begin with a custom \api{AbstractTableModel}, a
set of integer constants, a renderer switch, a selection listener, and perhaps
a preference key based on a header. Do not rewrite such a table all at once
unless it is very small. A staged migration keeps the screen usable while
identity becomes stronger.

The first stage is naming the row type. If the model currently stores arrays
or maps, introduce a row record or row class at the boundary where values are
already known. The \api{getValueAt} method can still switch on columns for a
while. The important improvement is that row values have names and types.
Tests can start constructing rows directly.

The second stage is naming columns. Introduce \api{ColumnKey} constants for
the existing columns without changing rendering. Replace persistence keys,
validation targets, and tests first, because those are the places where stable
identity pays off immediately. The table can still use the old model while
surrounding code stops depending on raw indexes.

The third stage is building descriptors. Create \api{ColumnDescriptor} and
\api{Column} objects that mirror the existing model. At this point it is often
clear which old column constants were really view positions and which were
business identities. Write a small test that traces a descriptor column from
key to value extractor. This test will protect the migration when the Swing
model is replaced.

The fourth stage is replacing the model with \api{ObservableTableModel}. Use
the same row data and descriptor. Keep renderers and selection policy as close
as possible to the old screen. The goal of this stage is not visual change; it
is to put descriptor-backed values under the table. If the screen changes
appearance at this stage, it will be harder to tell whether the descriptor or
the view changed.

The fifth stage is moving selection into presenter state. Install
\api{TableSelectionBinding}. Change commands to read selected row or selected
id values. Leave the visible table selection behavior unchanged. Tests should
assert selection across sorting if sorting exists. This stage often deletes
several small listeners that used to duplicate coordinate conversion.

The sixth stage is installing table state by stable ids. Start with an
\api{InMemoryTableStateStore} in tests, then move to preferences or the
application store. If old user preferences exist, write a migration from old
keys to new persistence ids. If old preferences used localized headers, this
is the time to fix that permanently.

The seventh stage is moving feedback into bindings. Header tooltips, cell
tooltips, validation markers, and help topics can be installed from descriptors
and problem state. Renderers become simpler because they display state rather
than discovering it. Tests become clearer because problem identity and column
identity are typed.

The eighth stage is optional generation. Once a descriptor pattern is stable,
decide whether \api{@SwingTable} can own the row declaration. Some legacy
tables have dynamic columns or specialized row models and should remain
handwritten. Others are ordinary record-backed tables and can move to
generated companions. Generation is a convenience after the contract is
understood, not a substitute for understanding the contract.

During migration, keep screenshots. A table is visual and dense. A screenshot
before and after model replacement can catch lost widths, poor headers, dark
theme problems, hidden columns, row height changes, and renderer regressions.
Pair screenshot review with targeted tests for descriptors and selection so a
visual difference can be interpreted rather than guessed.

The migration is complete when a reviewer can answer these questions without
searching for integer conventions: what is the row type, what are the column
keys, where are headers resolved, where is selection stored, where is table
state persisted, where are validation problems targeted, where are renderers
installed, and who closes table bindings? If those answers are visible, the
legacy table has crossed the important boundary.

\section{Testing, review, and failure patterns}

Table tests should be explicit about coordinate systems. A table state
assertion should use persisted column ids. A selection assertion should
distinguish view row, model row, row value, and row identity. A validation
assertion should target row identity and column key. A renderer test should
not pretend to prove validation semantics. A screenshot should prove the table
renders, not that a service query is correct.

\begin{lstlisting}[caption={Testing selection after sorting by row value.}]
tester.selectTableModelRow(CustomerComponents.TABLE, 1);
tester.sortBy(CustomerColumns.DISPLAY_NAME);

tester.assertSelectedRowValue(
        CustomerComponents.TABLE,
        row -> row.id().equals(CustomerId.of(42)));
\end{lstlisting}

The exact helper names may vary, but the policy does not. Tests should say
which coordinate system they are asserting. If the test wants current view
position, say view row. If it wants source position, say model row. If it wants
domain continuity, say row id. Ambiguous table tests often pass until the
first sorter, filter, or hidden column is added.

Review table code by tracing one column through the whole screen. Start with
the source record component or handwritten column definition. Find its
\api{ColumnKey}. Find the header resource key. Find the value extractor. Find
the renderer, if any. Find the validation target. Find the persisted column
state. Find export and help behavior if they exist. If each step uses the same
descriptor identity, the table is coherent. If one step switches to header
text or view index, that step is a likely future bug.

The first failure pattern is column-index sprawl. A model switch, renderer
switch, export switch, validation switch, and tooltip switch all use integers.
Recovery begins by creating column keys and descriptors. Do not attempt to fix
every switch at once if the table is large. Replace one column identity, then
move consumers to it.

The second failure pattern is selection by visible row. It works until sorting
or filtering arrives. Recovery begins by binding selection to model row or row
identity and by changing command handlers to read presenter selection state
instead of \api{JTable.getSelectedRow}. This makes master-detail screens much
easier to test.

The third failure pattern is persisted state by header. It works until
localization or copyediting. Recovery begins by assigning column persistence
ids and migrating old stored state if necessary. Do not change ids casually
after users have saved layouts.

The fourth failure pattern is renderer-owned behavior. A renderer computes
validation, loads icons from disk, formats with newly allocated formatters, or
records state in fields. Recovery begins by moving facts out of the renderer:
validation into problem sets, icons into caches, formatting policy into helper
objects, and row identity into descriptors.

The fifth failure pattern is hidden table lifecycle. A table state controller,
tooltip binding, selection listener, and model subscription are installed, but
no scope closes them. Recovery begins by listing every installed relationship
and adding a view-owned binding scope. A table is often the largest listener
cluster in a screen; its cleanup should be visible.

\section{Operational recipes}

The first recipe is a sorted master-detail table. Define a row type with a
durable id. Define a descriptor with stable column keys. Build an observable
row source. Create the \api{ObservableTableModel} on the EDT. Install a
\api{TableRowSorter} if sorting is local. Install \api{TableSelectionBinding}
after the sorter. Bind selected row identity into presenter state. Let the
detail panel observe selected identity or selected row, then load detail
through a task when necessary. Preserve selection across refresh by durable id,
not by old view row.

The second recipe is an editable table of immutable rows. Use editable
\api{Column} definitions whose updater returns a replacement row. Keep
parsing policy in editors or form-like cell models, and keep semantic
validation in problem models. On commit, let the table model replace the row
in the observable list. If a save command commits table edits, explicitly stop
or cancel the active cell editor before reading row state. Test both a valid
edit and an invalid intermediate edit. The test should prove that row identity
is preserved or deliberately changed according to policy.

The third recipe is a table with a column chooser. Give every hideable column
a stable persistence id and a localized resource name. Install
\api{TableStateController}. Install a visibility menu that talks to the
controller by persistence id. Do not remove columns from the model to hide
them. Do not store hidden state by header text. Test that hiding and restoring
a column preserves column order, width, and sort policy where appropriate.
Then change the header resource text and prove the stored visibility still
applies.

The fourth recipe is a validation-heavy table. Decide whether problems target
row object, row id, or model row. Use \api{ColumnKey} for the column side.
Compute problems outside renderers. Install \api{TableCellValidationBinding}
to decorate cells. Compose cell tooltips from problem, static help, and
formatted value. Add a problem summary if keyboard users need a stable place
to inspect errors. Test validation by problem identity, not by the exact
rendered red border or localized message.

The fifth recipe is a table with background refresh. Commands start tasks.
Tasks load rows away from the EDT. Results return through an explicit EDT
boundary or an EDT-adapted list. A generation counter suppresses stale
results. The row source applies precise changes when that effort is justified.
Selection is restored by durable id. Table state remains untouched by row
refresh. Busy state disables incompatible commands and may show an overlay.
Failure state says whether rows are stale, cleared, or still usable. None of
this belongs in a renderer or table model.

The sixth recipe is export. Decide whether export follows visible columns,
descriptor columns, or a separate export descriptor. Resolve export headers
through resources or export-specific labels. Read values through descriptor
columns or export value extractors. Convert selected view rows to row
identity or model rows before export. Do not scrape renderer components. Test
export after column reorder, after hiding a column, and after localization.
Those tests catch whether export is accidentally coupled to the current view.

The seventh recipe is row detail navigation. A double-click, Enter key, popup
command, and toolbar command should route through the same selected-row or
selected-id state. The table event should translate user gesture into command
execution, not directly open a service or dialog. This keeps keyboard and
mouse behavior consistent and lets command enablement explain why navigation
is unavailable. If double-click on a row should select it first, make that
policy explicit and test it with sorting active.

The eighth recipe is table help. Attach help topics to column descriptors.
Install header or cell help behavior that identifies the current column by
model column and descriptor key. Include row context only when help needs it.
Use F1 and explicit help commands rather than hover-only documentation. A help
request should contain a topic and structured context, not visible header
text. Test help dispatch with a fake help service; no browser is needed to
prove the boundary.

The ninth recipe is teardown. Put the table model, selection binding, state
controller, validation binding, tooltip binding, renderer installers, popup
listeners, and key bindings into the view's scope. Close the scope when the
view closes. Then mutate the row source and problem state in a test and prove
the old component is no longer updated. Close pending table-state saves or
flush them according to policy. A table is too listener-rich to rely on
garbage collection as cleanup.

These recipes are intentionally repetitive in one respect: every one starts
from typed identity and ends with scoped Swing relationships. That is the
Corusco table style. It does not prevent a screen from using ordinary Swing
features. It keeps those features attached to names that survive sorting,
localization, refresh, and view teardown.

\begin{tabularx}{0.96\linewidth}{@{}p{0.30\linewidth}X@{}}
\toprule
Scenario & Primary owner to inspect first\\
\midrule
Wrong value saved after sorting & Selection binding and row identity policy\\
Column widths lost after release & Table and column persistence ids, then migration\\
Tooltip disagrees with validation & Problem source and tooltip composition policy\\
Slow scrolling & Renderer allocation, row event volume, and sorter/filter owner\\
Export columns wrong after hiding & Export policy and descriptor/table-state boundary\\
Help opens generic page & Column help topics and help dispatch context\\
Old table updates after close & View scope and binding/controller cleanup\\
\bottomrule
\end{tabularx}

\section{Decision rules}

Use a generated \api{@SwingTable} descriptor when the row is naturally a record
and the column set is known at compile time. Use a handwritten descriptor when
columns are dynamic, domain-specific, or assembled from runtime configuration.
Use plain Swing only when the table is genuinely local, short-lived, and has
no durable column identity outside painting. This is not a hierarchy of virtue.
It is a choice about where the table's meaning can be stated most clearly.

Use Swing sorting when sorting is a view preference over a loaded local list.
Use a Corusco sorted list when sorted order is shared by several parts of the
screen. Use service-side sorting when the result is large, remote, paged, or
part of a query contract. The descriptor can mark a column sortable in all
three cases, but the sorting owner determines how selection, export, refresh,
and status text behave.

Use Swing filtering when the filter is local to one table and does not affect
other screen facts. Use a presenter-owned filtered list when commands,
summaries, export, or detail workflow need to know the filtered set. Use
service-side filtering when remote volume or security requires it. A filter
that users experience as "the current result set" usually belongs above the
\api{JTable}; a filter that users experience as "temporarily hide these rows in
this view" may belong in the sorter.

Use row-value selection when the row objects are stable enough for the
workflow. Use row-id selection when refresh, paging, or replacement rows are
common. Use model-row selection only for immediate table mechanics. Use view
row only inside Swing event handling and tests specifically about view
coordinates. This rule prevents many master-detail bugs because it makes the
selection's intended lifetime explicit.

Use table-state persistence when users can reasonably personalize the table:
widths, order, visibility, or sort. Do not persist state for tiny static
tables unless the application benefits from it. Once state is persisted, ids
become compatibility-sensitive. A throwaway table can rename columns freely; a
persisted business table needs migration discipline.

Use validation bindings when problems are computed outside rendering and need
to be shown in cells. Use custom renderers when values need special visual
presentation. Use both when a cell has both display policy and problem state.
Do not make the renderer compute the problem merely because the problem is
drawn there. Rendering is presentation; validation is model or presenter
policy.

Use descriptor-based export when the exported data has semantic meaning. Use
screenshot or copy helpers when the user explicitly wants visible text in
visible order. The first path is data export; the second is view capture. They
may share resource formatting, but they are different workflows. Confusing
them leads to exported icons, hidden columns, or localized headers in places
where stable schema was needed.

Use optimized renderers after measurement or when the workload is obviously
dense. Do not replace every default renderer preemptively. Swing's default
renderers are competent for many tables and cooperate with Look-and-Feel
defaults. Premature custom renderers often break selection colors,
accessibility, dark themes, or performance by accident. A descriptor-backed
table can begin plain and become specialized column by column.

These rules should be written into code review. When a table change arrives,
ask which owner was chosen and whether that owner matches the workflow. A
short answer is acceptable: "local sorter because this is a small loaded
result", "service sort because the query is paged", "row id selection because
refresh replaces row objects", "handwritten descriptor because columns come
from a user-defined report." The answer matters more than the particular
class imported.

A second review pass should ask whether the table can explain itself without
the screen running. The descriptor should name columns. The state object should
name persisted ids. The selection model should name its coordinate system. The
problem target should name row and column meaning. The example should name the
fixture data, and the screenshot should show the state the prose discusses. If
one of these facts exists only as a current Swing index, a localized header, or
a renderer side effect, the table is borrowing meaning from pixels. That can
work for a demo, but it is too weak for a table users personalize, export,
validate, or revisit after a restart.

\section{Exercises and checklist}

\begin{exercisebox}
\begin{enumerate}
\item Replace one column-index constant with a \api{ColumnKey} and update the
renderer, tooltip, and test that used the old index.
\item Add selection binding to a sorted table and verify that model-row state
remains correct after sorting.
\item Persist column widths by descriptor id, not header text, and change the
visible header to prove state survives.
\item Add a validation problem to a row and column key, then display it through
a cell validation binding.
\item Hide a column and restore it through table state rather than rebuilding
the model.
\item Profile a renderer that formats dates. Move formatter allocation out of
the rendering path.
\item Add a table-state migration for a real column split and document why the
old id cannot remain.
\item Add a table export path that uses descriptors rather than renderer text.
\item Close the table view scope and prove that later row-source changes do not
notify the old table model.
\item Compare a Swing \api{RowFilter} with a presenter-owned filtered list and
decide which owner fits the screen.
\end{enumerate}
\end{exercisebox}

\begin{center}
\textbf{Corusco table checklist.}
\end{center}

\begin{itemize}
\item Use descriptors for stable table meaning.
\item Keep view indexes out of persistence, validation identity, export, help,
and tests.
\item Adapt off-EDT row sources before Swing consumes them.
\item Bind selection by model row, row value, or row id when sorting and
filtering are active.
\item Persist table state through stable table and column ids.
\item Treat hidden columns as descriptor columns with invisible presentation,
not as deleted model columns.
\item Use migration only for real table-schema changes.
\item Compose tooltips and validation rather than hiding logic in renderers.
\item Keep renderers stateless, theme-aware, and cheap.
\item Commit editable cells deliberately before command execution.
\item Close table models, controllers, renderers with installed state, and
bindings with the view lifecycle.
\item Test table behavior by naming the coordinate system under test.
\end{itemize}

The point of a Corusco table is not to make \api{JTable} disappear. Swing's
table component remains a serious, mutable, highly capable widget. The point is
to stop asking the widget to be the only place where table meaning lives. Once
row identity, column identity, resources, validation, state, selection, and
lifecycle have owners, a \api{JTable} can do what it does well: present data,
let users interact with it, and repaint quickly.
